\documentstyle[% 


righttag% 


,11pt% 


,amssymb% 


,russian%               remove in Eng. 


,izvru%                 Set izven% in Eng. 


]{amsart} 


 


% 


 


\newcommand{\const}{\operatorname{const}} 


\newcommand{\Div}{\operatorname{div}} 


\newcommand{\Ric}{\operatorname{Ric}} 


\newcommand{\id}{\operatorname{id}} 


\newcommand{\tts}[1]{} 


\numberwithin{equation}{section} 


 


\theoremstyle{plain} 


\theorembodyfont{\it} 


\newtheorem{thms}{\hskip\parindent Теорема}[section] 


\newtheorem{lems}{\hskip\parindent Лемма}[section] 


 


\theoremstyle{definition} 


\theorembodyfont{\rm} 


\newtheorem{cors}{\hskip\parindent Следствие}[section] 


\newtheorem{notes}{\hskip\parindent Замечание}[section]


 


%\hoffset=-15mm%                  For Eng. 


\setcounter{page}{69} 


\god{2004} 


\nomer{5~(504)} %                        Remove in Eng. 


%\nomer{Vol.\,48, No.\,5}%               Set in Eng. 


%\str{\pageref{begin}--\pageref{end}}%  Set in Eng. 


\udk{514.76} 


 


\author{\it М.В.\,Смольникова, С.Е.\,Степанов, И.Г.\,Шандра} 


% NB:   ^^ remove in Eng. 


\title{Инфинитезимальные гармонические преобразования} 


\thanks{Работа выполнена при финансовой поддержке Российского фонда 


фундаментальных исследований (проект \No\,03-01-00028).} 


 


\date{} 


%\pagestyle{myheadings}%         Set in Eng. 


\pagestyle{plain} %              Remove in Eng. 


\begin{document} 


%\thispagestyle{plain}%          Set in Eng. 


\maketitle 


\label{begin} 


 


 


 


 


\section*{Введение} 


 


Несмотря на обилие публикаций по гармоническим отображениям многообразий 


(см. обзоры [1]--[3]) на сегодняшний день нет определения 


инфинитезимального 


гармонического преобразования, не говоря уже о теории групп 


таких преобразований. Это резко контрастирует с теориями изометрических, 


конформных, проективных и других видов отображений и соответствующих им 


инфинитезимальных преобразований (см., напр., [4]--[9]). 


 


Данной работой авторы намерены восполнить этот пробел. В статье 


изучаются инфинитезимальные гармонические преобразования и локальные 


группы инфинитезимальных гармонических преобразований римановых и 


келеровых многообразий. 


 


Содержание первых трех параграфов данной статьи было анонсировано в 


[10], а четвертого --- в [11]. 


 


 


\section{Гармонические отображения римановых многообразий} 


 


1.1. Рассмотрим римановы многообразия $(M,g)$ и $(\ov M,\ov g)$ 


размерностей $m$ 


и $n$ соответственно и гладкое отображение $f:M\to\ov M$. Обозначим через 


$\{x^1,\dotsc,x^m\}$ локальные координаты в окрестности $U$ точки $x\in M$, 


$\{\ov x^1,\dotsc,\ov x^n\}$ --- локальные координаты в окрестности $\ov U\supset f(U)$ 


точки 


$\ov x=f(x)\in\ov M$, $\Gamma^k_{ij}$, $i,j,k=1,\dotsc,m$, и 


$\ov\Gamma^\alpha_{\beta\gamma}$, $\alpha,\beta,\gamma=1,\dotsc,n$, 


--- символы Кристофеля связностей Леви-Чивита $\nabla$ и $\ov\nabla$ римановых 


многообразий 


$(M,g)$ и $(\ov M,\ov g)$ соответственно. Для того чтобы гладкое 


отображение 


$f:M\to\ov M$ римановых многообразий $(M,g)$ и $(\ov M,\ov g)$ было {\it 


гармоническим} 


(см., напр., [1]--[3]), необходимо и достаточно, чтобы выполнялись 


уравнения Эйлера--Лагранжа 


\begin{equation} 


g^{ij}(\partial_i\partial_j f^\alpha-\Gamma^k_{ij} 


\partial_k f^\alpha+\ov\Gamma^\alpha_{\beta\gamma} 


\partial_i f^\beta\partial_j f^\gamma)=0 


\end{equation} 


для контрвариантных компонент $g^{ij}$ метрического тензора $g$ и 


$\partial_i=\frac{\partial}{\partial x^i}$. 


 


\begin{notes} 


В случае евклидова пространства уравнения (1.1) превращаются 


в систему уравнений Лапласа--Бельтрами $\Delta 


f^\alpha=\sum\limits^m_{j=1} \partial^2_j f^\alpha=0$, что и объясняет 


происхождение термина {\it гармоническое отображение}. 


\end{notes} 


 


1.2. Если $\dim M=\dim\ov M=m$ и отображение $f:M\to\ov M$ является 


диффеоморфизмом, 


то локально отображение $f$ будет осуществляться ([12], с.\,67) 


по принципу равенства координат $\ov x^1=x^1,\dotsc,\ov x^m=x^m$ 


соответствующих точек 


$\ov x=f(x)$. В этом случае уравнения Эйлера--Лагранжа (1.1) предстают в 


следующем виде: 


\begin{equation} 


g^{ij}T^k_{ij}=0, 


\end{equation} 


где $T^k_{ij}=\ov\Gamma^k_{ij}-\Gamma^k_{ij}$ --- компоненты {\it тензора 


деформации} $T=\ov\nabla-\nabla$ связности Леви-Чивита 


$\nabla$ многообразия $(M,g)$ при отображении $f$ на многообразие 


$(\ov M,\ov g)$ ([4], с.\,162; [6], с.\,71). Отметим, что для 


произвольного 


диффеоморфизма $f:M\to\ov M$ тензор деформации $T=\ov\nabla-\nabla$ 


рассматривают (см. там 


же) как гладкое сечение расслоения $TM\otimes S^2M$. В случае же 


гармонического 


диффеоморфизма тензор деформации $T=\ov\nabla-\nabla$ в силу условия (1.2) 


становится 


бесследовым симметрическим тензорным полем --- сечением расслоения 


$f^{-1}TM'\otimes S^2_0M$. 


Таким образом, доказана 


 


 


\begin{lems} 


Для того чтобы диффеоморфизм $f:M\to\ov M$ риманова многообразия $(M,g)$ 


на риманово многообразие $(\ov M,\ov g)$ был гармоническим, необходимо и 


достаточно, чтобы тензор деформации $T=\ov\nabla-\nabla$ связности 


Леви-Чивита 


$\nabla$ многообразия $(M,g)$ при отображении $f$ на многообразие $(\ov 


M,\ov g)$ 


был сечением тензорного расслоения $f^{-1}TM'\otimes S^2_0M$. 


\end{lems} 


 


 


\section{Определение и примеры 


инфинитезимального гармонического преобразования} 


 


2.1. Произвольное векторное поле $\xi\in C^\infty TM$ порождает в 


окрестности $U$ 


каждой точки $x$ риманова многообразия $(M,g)$ локальную 1-параметрическую 


группу инфинитезимальных преобразований $\varphi_t:U\to M$, задаваемых 


формулами $\varphi_t(x)=\ov x^k=x^k+t\xi^k$ для локальной системы 


координат 


$\{x^1,\dotsc,x^m\}$ в окрестности $U$, параметра 


$t\in(-\varepsilon,+\varepsilon)\subset \bold R$ и $\xi=\xi^k\partial_k$ 


([13], 


с.\,21--23; [14], с.\,39--41). На этом основании векторное поле $\xi\in 


C^\infty TM$ называют 


еще {\it инфинитезимальным преобразованием в} $(M,g)$. 


 


Локальные компоненты $g_{ij}$ метрического тензора $g$ и символы 


Кристофеля 


$\Gamma^k_{ij}$ связности Леви-Чивита $\nabla$ в результате 


инфинитезимального преобразования 


предстанут в следующем виде ([14], сс.\,40, 41): 


\begin{equation} 


\ov g_{ij}=g_{ij}+t(L_\xi\,g_{ij}),\quad 


\ov\Gamma^k_{ij}=\Gamma^k_{ij}+t(L_\xi\,\Gamma^k_{ij}). 


\end{equation} 


Здесь 


\begin{equation} 


L_\xi\,g_{ij}=\nabla_i\xi_j+\nabla_j\xi_i;\quad 


L_\xi\,\Gamma^k_{ij}=\nabla_i\nabla_j\xi^k-R^k_{ijl}\xi^l, 


\end{equation} 


где $R^k_{ijl}$ --- компоненты тензора кривизны $R$ связности $\nabla$, 


являются выражениями 


производных Ли ([13], с.\,37) компонент $g_{ij}$ тензора $g$ и символов 


Кристофеля 


$\Gamma^k_{ij}$ по отношению к векторному полю $\xi$. 


 


В соответствии с общей теорией ([9], с.\,9; [13], с.\,19) гармонический 


диффеоморфизм $f$ риманова многообразия $(M,g)$ на себя назовем {\it 


гармоническим 


преобразованием} в многообразии $(M,g)$. Аналогично, векторное поле $\xi$ 


назовем {\it инфинитезимальным гармоническим преобразованием} в $(M,g)$, 


если 


локальная однопараметрическая группа инфинитезимальных преобразований, 


порожденная векторным полем $\xi$ в окрестности $U$ каждой точки $x$ 


многообразия 


$(M,g)$, состоит из инфинитезимальных гармонических преобразований. 


 


Согласно равенствам (2.1) заключаем, что определяющими для 


инфинитезимального 


гармонического преобразования $\xi$ будут уравнения 


\begin{equation} 


g^{ij}(L_\xi\,\Gamma^k_{ij})=0. 


\end{equation} 


 


\begin{thms} 


Для того чтобы векторное поле $\xi\in C^\infty TM$ было инфинитезимальным 


гармоническим преобразованием в $(M,g)$, необходимо и достаточно, 


чтобы его компоненты удовлетворяли дифференциальным уравнениям 


$\Delta\xi^k=2R^k_j\xi^j$. 


\end{thms} 


 


\begin{pf} 


Действие оператора Лапласа--Бельтрами $\Delta$ на произвольное 


векторное поле $\xi$ выражается ([9], с.\,203) следующей формулой: 


$\Delta\xi^k=-g^{ij}\nabla_i\nabla_j\xi^k+ R^k_j\xi^j$, где 


$R_{ij}=g_{ik}R^k_j$ --- компоненты тензора Риччи $Ric$  %%% put \Ric in Eng.


многообразия $(M,g)$. С другой стороны, на основании равенств (2.2) 


уравнениям 


(2.3) придадим следующий равносильный им вид: 


\begin{equation} 


g^{ij}\nabla_i\nabla_j\xi^k+R^k_j\xi^j=0. 


\end{equation} 


Из этих двух систем уравнений выводим $\Delta\xi^k=2R^k_j\xi^j$. 


\end{pf} 


 


\begin{notes} 


В случае евклидова пространства уравнения (2.4) предстают в виде 


системы уравнений Лапласа--Бельтрами $\Delta\xi^k= 


\sum\limits^m_{j=1}\partial^2_j\xi^k=0$, что объясняет происхождение 


термина {\it инфинитезимальное гармоническое преобразование}. 


\end{notes} 


 


2.2. Рассмотрим примеры инфинитезимальных гармонических преобразований. 


Первым примером будет инфинитезимальное конформное преобразование 


в $(M,g)$, под которым понимают ([13], с.\,284; [14], с.\,47) векторное 


поле $\xi$, подчиняющееся условию 


\begin{equation} 


L_\xi\,g_{ij}=\nabla_i\xi_j+\nabla_j\xi_i=\frac{2}{m} 


(\nabla_k\xi^k)g_{ij}. 


\end{equation} 


При этом поле $\xi$ называют {\it инфинитезимальной гомотетией} 


(соответственно {\it инфинитезимальной изометрией}), если в дополнение к 


(2.5) выполняется 


условие $\nabla_k\xi^k=\const$ (соответственно $\nabla_k\xi^k=0$). 


Локальная однопараметрическая 


группа инфинитезимальных преобразований, порождаемая полем $\xi$, будет 


группой локальных конформных преобразований (соответственно группой 


локальных 


гомотетий или изометрий) тогда и только тогда, когда $\xi$ --- 


инфинитезимальное 


конформное преобразование (соответственно инфинитезимальная 


гомотетия или инфинитезимальная изометрия). 


 


В случае конформного преобразования $\xi$ непосредственные расчеты 


показывают 


([5], гл.\,2, \S\,5), что $\Delta\xi_k= 2R_{kj}\xi^j+ m^{-1}(m-2)\nabla_k 


(\nabla_j\xi^j)$. А потому каждое инфинитезимальное 


конформное преобразование двумерного риманова 


многообразия $(M,g)$ является гармоническим преобразованием. Для 


размерностей 


$n\ge3$ это возможно только в случае, когда данное инфинитезимальное 


преобразование является либо гомотетией, либо изометрией. 


 


Обратимся теперь к келерову многообразию $(M,g,J)$, которое является 


римановым многообразием $(M,g)$ четной размерности $m=2n$ с почти 


комплексной 


структурой $J$, которая в каждой точке $x\in M$ является эндоморфизмом 


касательного пространства $T_xM$ таким, что $J^2=-\id$, и при этом 


$g(J,J)=g$ и $\nabla J=0$. Векторное поле $\xi$, задающее 


инфинитезимальный автоморфизм 


почти комплексной структуры келерова многообразия $L_\xi\,J=0$, называется 


(вещественным) {\it голоморфным векторным полем} ([15], с.\,120). 


 


К.\,Яно доказал ([5], гл.\,4, \S\,6, теорема 6.5), что уравнения (2.4) являются 


следствиями уравнения $L_\xi\,J=0$ и что в случае компактного келерова 


многообразия $(M,g,J)$ верно и обратное утверждение. Таким образом, 


справедлива 


 


\begin{thms} 


Каждое $($вещественное$)$ голоморфное векторное поле на келеровом 


многообразии $(M,g,J)$ является инфинитезимальным гармоническим 


преобразованием. В случае компактного многообразия $(M,g,J)$ любое 


инфинитезимальное гармоническое преобразование в $(M,g,J)$ является 


$($вещественным$)$ голоморфным векторным полем. 


\end{thms} 


 


 


\section{Свойства инфинитезимального гармонического преобразования} 


 


 


Выше было установлено, что инфинитезимальное изометрическое 


преобразование в римановом многообразии $(M,g)$ является гармоническим. 


Докажем теперь, что справедлива 


 


\begin{thms} 


Для того чтобы инфинитезимальное гармоническое преобразование $\xi$ 


в римановом многообразии $(M,g)$ было изометрическим, достаточно, 


чтобы многообразие $(M,g)$ было компактным, а поле $\xi$ --- 


соленоидальным. 


\end{thms} 


 


\begin{pf} 


Напомним ([9], с.\,63), что равенства $\Delta\xi^k=2R^k_j\xi^j$ и 


$\nabla_j\xi^j=0$ являются необходимым и достаточным условием для того, 


чтобы $\xi$ 


было инфинитезимальной изометрией компактного многообразия $(M,g)$. 


Поэтому для доказательства достаточно обратиться к теореме 2.1. 


\end{pf} 


 


В [16] было введено понятие {\it гармонического тензорного поля} $\varphi$ 


на 


римановом многообразии $(M,g)$ как сечения расслоения $S^2M$, локальные 


компоненты $\varphi_{jk}$ которого удовлетворяют дифференциальным 


уравнениям 


$g^{ij}\nabla_i\varphi_{jk}=0$. Охарактеризуем инфинитезимальное 


гармоническое преобразование 


$\xi$ в многообразии $(M,g)$ в терминах гармонического тензорного поля 


$\varphi$. 


 


\begin{thms} 


Векторное поле $\xi$ будет инфинитезимальным гармоническим 


преобразованием в римановом многообразии $(M,g)$ тогда и только тогда, 


когда 


тензорное поле $\varphi=L_\xi g-(\Div\xi)g$ является гармоническим на 


$(M,g)$. 


\end{thms} 


 


\begin{pf} 


Для тензорного поля $\varphi=L_\xi g-(\Div\xi)g$ выпишем 


три группы очевидных равенств: 


\begin{align*} 


\nabla_k\nabla_j\xi_i+\nabla_k\nabla_i\xi_j- 


\nabla_k(\nabla_l\xi^l)g_{ij}&=\nabla_k\varphi_{ij},\\ 


% 


\nabla_j\nabla_k\xi_i+\nabla_j\nabla_i\xi_k- 


\nabla_j(\nabla_l\xi^l)g_{ik}&=\nabla_j\varphi_{ik},\\ 


%


\nabla_i\nabla_j\xi_k+\nabla_i\nabla_k\xi_j- 


\nabla_i(\nabla_l\xi^l)g_{kj}&=\nabla_i\varphi_{jk}.


\end{align*} 


Затем сложим почленно уравнения первой и второй групп и вычтем из 


полученного результата уравнения третьей группы, в результате получим 


\begin{multline*} 


\nabla_k\nabla_j\xi_i+\nabla_j\nabla_k\xi_i-\xi_l R^l_{jki} 


-\xi_lR^l_{kji}-\nabla_k(\nabla_l\xi^l)g_{ij}- \\ 


% 


-\nabla_j(\nabla_l\xi^l)g_{ik}+\nabla_i(\nabla_l\xi^l)g_{jk}= 


\nabla_k\varphi_{ij}+\nabla_j\varphi_{ki}-\nabla_i\varphi_{jk}. 


\end{multline*} 


Последние уравнения после свертки их левой и правой частей с 


контрвариантными компонентами $g^{ki}$ метрического тензора $g$ 


примут следующий вид: 


$$ 


g^{kj}\nabla_k\nabla_j\xi_i+\xi_lR^l_i= 


g^{jk}\nabla_j\varphi_{ki}. 


$$ 


Теперь очевидно, что условия гармоничности инфинитезимального 


преобразования 


$\xi$ и гармоничности тензорного поля $\varphi=L_\xi g-(\Div\xi)g$ 


равносильны. 


\end{pf} 


 


Хорошо известны ``теоремы исчезновения'' для инфинитезимальных конформных 


и изометрических преобразований ([13], с.\,235; [15], с.\,48). Докажем 


аналогичную теорему и для инфинитезимальных гармонических преобразований. 


 


\begin{thms} 


Компактное риманово многообразие $(M,g)$ не допускает отличных 


от нуля инфинитезимальных гармонических преобразований, если кривизна 


Риччи этого многообразия отрицательна. 


\end{thms} 


 


\begin{pf} 


Применим оператор Лапласа--Бельтрами $\Delta$ к функции 


$F=g_{ij}\xi^i\xi^j$ для произвольного инфинитезимального гармонического 


преобразования 


$\xi$ многообразия $(M,g)$. Найдем, что 


\begin{equation} 


\Delta F=g^{ij}\nabla_i\nabla_j(g_{ij}\xi^i\xi^j)= 


2[-R_{ij}\xi^i\xi^j+g^{kl}g_{ij} 


(\nabla_k\xi^i)(\nabla_l\xi^j)]. 


\end{equation} 


Далее, на основании предположения о кривизне Риччи многообразия $(M,g)$ 


заключаем, что $\Delta F\ge0$. Для завершения доказательства достаточно 


воспользоваться 


леммой Хопфа ([17], с.\,308; [14], сс.\,29, 32--34), в силу которой из (3.1) 


следует $\xi=0$. 


\end{pf} 


 


 


\begin{thms} 


Инфинитезимальное гармоническое преобразование $\xi$ в компактном 


келеровом многообразии $(M,g,J)$ постоянной скалярной кривизны 


разлагается в сумму $\xi=\xi'+J\xi''$, где $\xi'$ и $\xi''$ являются 


инфинитезимальными 


изометриями в $(M,g,J)$. 


\end{thms} 


 


\begin{pf} 


Согласно теореме 2.3 каждое инфинитезимальное гармоническое 


преобразование $\xi$ компактного келерова многообразия $(M,g,J)$ является 


голоморфным (вещественным) преобразованием в $(M,g,J)$. С другой 


стороны, по теореме А.\,Лихнеровича ([5], с.\,92--98) подобное 


векторное поле $\xi$ на компактном келеровом многообразии $(M,g,J)$ 


постоянной 


скалярной кривизны разлагается в сумму $\xi=\xi'+J\xi''$, где $\xi'$ и 


$\xi''$ суть инфинитезимальные изометрии в $(M,g,J)$.


\end{pf} 


 


 


\section{Векторное пространство инфинитезимальных 


гармонических преобразований} 


 


4.1. Рассмотрим пространство сечений $C^\infty S^pM$ векторного расслоения 


$S^pM$ ковариантных симметрических $p$-тензоров на $n$-мерном компактном 


ориентированном римановом многообразии $(M,g)$. В соответствии с общей 


теорией ([5], гл.\,4, \S\,1) полагаем 


\begin{equation} 


\langle\varphi,\ov\varphi \rangle=\int_M\frac{1}{p!} 


g(\varphi,\varphi') 


\end{equation} 


для $g(\varphi,\varphi')= \varphi^{i_1\dotsc i_p} \varphi'_{i_1\dotsc 


i_p}= g^{i_1k_1}\dotsc g^{i_p k_p} \varphi_{k_1\dotsc k_p} 


\varphi'_{i_1\dotsc i_p}$ и локальных компонент 


$\varphi'_{i_1\dotsc i_p}$ и $\varphi_{k_1\dotsc k_p}$ тензорных полей 


$\varphi,\varphi'\in C^\infty S^pM$. 


 


Построим векторное поле $X$ из локальных компонент $\varphi_{i_1\dotsc 


i_p}$ и 


$\psi_{ji_1\dotsc i_p}$ тензорных полей $\varphi\in C^\infty S^pM$ и 


$\psi\in C^\infty S^{p+1}M$, используя 


равенства $X^j=\varphi_{i_1\dotsc i_p}\psi^{ji_1\dotsc i_p}$. Далее, на 


основании теоремы Грина ([13], 


с.\,259) $\int\limits_M \nabla_i X^i dV=0$, и используя определенное 


глобально скалярное произведение 


(4.1), выводим равенство 


\begin{equation} 


\langle \delta^*\varphi,\psi\rangle=\langle 


\varphi,\delta\psi\rangle. 


\end{equation} 


Здесь $\delta^*:C^\infty S^pM\to C^\infty S^{p+1}M$ и $\delta:C^\infty 


S^{p+1}M\to C^\infty S^pM$ суть дифференциальные операторы 


первого порядка, определяемые следующими формулами 


(ср. с [15], с.\,54--55): 


$$ 


(\delta^*\varphi)_{i_0i_1\dotsc i_p}=\nabla_{i_0} 


\varphi_{i_1\dotsc i_{p-1}i_p}+\dotsb+ 


\nabla_{i_p}\varphi_{i_0i_1\dotsc i_{p-1}},\quad 


(\delta\psi)_{i_1\dotsc i_p}=-g^{ji_0} 


\nabla_j\psi_{i_0i_1\dotsc i_p}. 


$$ 


Равенство (4.2) свидетельствует о формальной сопряженности операторов 


$\delta^*$ и $\delta$ ([15], с.\,626; [18], с.\,69--70). 


\medskip 


 


4.2. Определим на пространстве $C^\infty S^pM$ дифференциальный оператор 


второго порядка $\square:C^\infty S^pM\to C^\infty S^pM$ следующей 


формулой: $\square=\delta\delta^*-\delta^*\delta$. 


Непосредственные расчеты показывают, что оператор $\square$ является 


самосопряженным, т.\,е. $\langle \square\varphi, \varphi'\rangle= 


\langle\varphi,\square\varphi' \rangle$. 


 


Полагая далее $\theta\in T^*_xM-\{0\}$ и $\varphi_x\in S^p_x M$, найдем 


символ $\sigma(\square)$ дифференциального 


оператора $\square$ ([15], с.\,627--628; [18], сс.\,64, 79, 87). Имеем 


$$ 


\sigma(\square)(\theta,x)\varphi_x=-i_{\theta^{\#}} 


(\theta\circ\varphi_x)+\theta\circ i_{\theta^{\#}}\varphi_x= 


-g(\theta,\theta)\varphi_x 


$$ 


для симметрического тензорного умножения $\circ$. Последнее равенство 


означает, 


что самосопряженный дифференциальный оператор второго порядка $\square$ 


является 


лапласианом ([15], с.\,77), а его ядро --- конечномерным векторным 


пространством ([15], с.\,631--632; [18], с.\,178). 


 


Поскольку 


\begin{align*} 


(\delta\delta^*\varphi)_{i_1i_2\dotsc i_{p-1}i_p}&= 


-g^{ji}(\nabla_i\nabla_j 


\varphi_{i_1i_2\dotsc i_{p-1}i_p}+\nabla_i\nabla_{i_1} 


\varphi_{i_2\dotsc i_{p-1}i_p j}+\dotsb+\nabla_i\nabla_{i_p} 


\varphi_{ji_1\dotsc i_{p-1}});\\ 


% 


(\delta^*\delta\varphi)_{i_1i_2i_3\dotsc i_{p-1}i_p}&= 


-g^{ji}(\nabla_{i_1}\nabla_i\varphi_{ji_2\dotsc i_p}+ 


\nabla_{i_2}\nabla_i\varphi_{ji_3\dotsc i_p i_1}+\dotsb+ 


\nabla_{i_p}\nabla_i\varphi_{ji_1\dotsc i_{p-1}}), 


\end{align*} 


то используя тождество Риччи ([4], с.\,42--43) найдем результат действия 


лапласиана $\square=\delta\delta^*-\delta^*\delta$ на симметрическое 


тензорное поле $\varphi$: 


\begin{multline} 


(\square\varphi)_{i_1\dotsc i_p}=-g^{ij}\nabla_i\nabla_j 


\varphi_{i_1\dotsc i_p}-\sum^p_{a=1}R^k_{i_a} 


\varphi_{i_1\dotsc i_{a-1}k i_{a+1}\dotsc i_p}- \\


-g^{kj}g^{lm}\sum^p\begin{Sb} a,b=1,\\a<b\end{Sb} 


R_{i_a j i_b m}\varphi_{i_1\dotsc i_{a-1}k 


i_{a+1}\dotsc i_{b-1}j i_{b+1}\dotsc i_p}. 


\end{multline} 


 


4.3. Для векторного поля $\xi$ введем дифференциальную 


форму $\theta$ такую, что $\theta^{\#}=\xi$ ([15], с.\,47). Имеем 


\begin{gather*} 


(\delta^*\theta)_{ij}=\nabla_i\xi_j+\nabla_j\xi_i,\quad 


(\delta\theta)=-g^{ij}\nabla_i\xi_i,\quad 


(\delta^*\delta\theta)_j=-g^{ik}\nabla_j\nabla_k\xi_i,\\ 


% 


(\delta\delta^*\theta)_j=-g^{ik}(\nabla_k\nabla_i\xi_j+ 


\nabla_k\nabla_j\xi_i). 


\end{gather*} 


Теперь, используя тождество Риччи ([4], с.\,42--43) найдем результат 


действия дифференциального оператора второго порядка 


$\square=\delta\delta^*-\delta^*\delta$ на дифференциальную 


форму $\theta$ (ср. с (4.3)) 


\begin{equation} 


(\square\theta)_j=[(\delta\delta^*-\delta^*\delta) 


\theta]_j=-g^{ik}(\nabla_k\nabla_i\xi_j- 


\xi_l R^l_{ikj})=-g^{ik}\nabla_k\nabla_i\xi_j+ 


R^l_j\xi_l. 


\end{equation} 


На основании равенства (4.4) заключаем, что справедлива 


 


\begin{thms} 


Инфинитезимальные гармонические преобразования в римановом 


многообразии $(M,g)$ и только они составляют ядро лапласиана 


$\square=\delta\delta^*-\delta^*\delta$. 


\end{thms} 


 


Нетрудно усмотреть, что инфинитезимальные гармонические преобразования 


в римановом многообразии $(M,g)$ образуют $\bold R$-модуль, который 


обозначим 


через $\cal H(M,\bold R)$. Поскольку оператор $\square$ является 


лапласианом, то справедливо 


 


\begin{cors} 


На компактном римановом многообразии $(M,g)$ \ $\bold R$-модуль 


$\cal H(M,\bold R)$ инфинитезимальных гармонических преобразований имеет 


конечную размерность. 


\end{cors} 


 


В заключение заметим, что на компактном многообразии $(M,g)$ с 


отрицательной 


кривизной Риччи $\dim\cal H(M,\bold R)=0$ согласно теореме 3.3. 


 


 


\begin{thebibliography}{99} 


 


\bibitem{1} 


Eells J., Lemaire L. {\em A report on harmonic maps} // Bull. London Math. 


Soc. -- 1978. -- V.\,10. -- \No\,1. -- P.\,1--68. 


\bibitem{2} 


Eells J., Lemaire L. {\em Another report on harmonic maps} // Bull. London 


Math. Soc. -- 1988. -- V.\,20. -- P.\,385--584. 


\bibitem{3} 


Давидов Й., Сергеев А.Г. {\em Твисторные пространства и гармонические 


отображения} // УМН. -- 1993. -- Т.\,48. -- \No\,3. -- С.\,3--96. 


\bibitem{4} 


Эйзенхарт Л.П. {\em Риманова геометрия}. -- М.: Ин. лит., 1948. -- 316 с. 


\bibitem{5} 


Yano К. {\em Differential geometry on complex and almost complex spaces}. 


-- Oxford: Pergamon Press, 1965. -- 326 p. 


\bibitem{6} 


Синюков Н.С. {\em Геодезические отображения римановых пространств}. -- М.: 


Наука, 1979. -- 255~с. 


\bibitem{7} 


Аминова А.В. {\em Группы преобразований римановых многообразий} // Итоги 


науки и техн. Пробл. геометрии. -- ВИНИТИ, 1990. -- Т.\,22. 


-- С.\,97--165. 


\bibitem{8} 


Yano K. {\em The theory of Lie derivatives and its applications}. -- 


Amsterdam: North Holland, 1957. -- 299 p. 


\bibitem{9} 


Кобаяси Ш. {\em Группы преобразований в дифференциальной геометрии}. -- 


М.: Наука, 1986. -- 224 с. 


\bibitem{10} 


Stepanov S.E. 


{\em The classification of harmonic diffeomorphisms} // Abstracts of the 


5${}^{\mathrm{th}}$ International Conf. on Geom. and Appl. August 


24--29, 2001, Varna. -- 


Sofia: Union of Bulgarian Mathematicians. -- 2001. -- P.\,55. 


\bibitem{11} 


Smolnikova M.V. {\em On global geometry harmonic symmetric bilinear 


differential 


forms} // Тез. докл. международн. конф. по дифференциальным 


уравнениям и динамическим системам, 21--26 августа 2000 г., Суздаль. 


-- Владимир: Изд-во Владимирск. ун-та, 2000. -- С.\,87-88. 


\bibitem{12} 


Нарасимхан Р. {\em Анализ на действительных и комплексных многообразиях}. 


-- М.: Мир, 1971. -- 232 с. 


\bibitem{13} 


Кобаяси Ш., Номидзу К. {\em Основы дифференциальной геометрии}. Т.\,1. 


-- М.: Наука, 1981. -- 344 с. 


\bibitem{14} 


Яно К., Бохнер С. {\em Кривизна и числа Бетти}. -- М.: Ин. лит., 1957. 


-- 152 с. 


\bibitem{15} 


Бессе А. {\em Многообразия Эйнштейна}. -- М.: Мир, 1990. -- 703 с. 


\bibitem{16} 


Garcia-Rio E., Vanhecke L., Vazquez-Abal E. {\em Harmonic endomorphism 


fields} // Illinois Journal of Mathematics. -- 1997. -- V.\,41. 


-- \No\,1. -- P.\,23--30. 


\bibitem{17} 


Кобаяси Ш., Номидзу К. {\em Основы дифференциальной геометрии}. 


Т.\,2. -- М.: Наука, 1981. -- 414 с. 


\bibitem{18} 


Пале Р. {\em Семинар по теореме Атьи--Зингера об индексе}. 


-- М.: Мир, 1970. -- 359 с. 


 


\end{thebibliography} 


 


\vskip 2cm 


 


\post{\it Владимирский государственный}{\it Поступила} 


\post{\it педагогический университет}{20.02.2002} 


\smallskip 


\post{\it Московская финансовая академия}{} 


 


\label{end} 


\end{document} 


